\documentclass[conference]{IEEEtran}
\IEEEoverridecommandlockouts
% The preceding line is only needed to identify funding in the first footnote. If that is unneeded, please comment it out.
\makeatletter
\newcommand{\linebreakand}{%
  \end{@IEEEauthorhalign}
  \hfill\mbox{}\par
  \mbox{}\hfill\begin{@IEEEauthorhalign}
}
\makeatother

\usepackage{cite}
\usepackage{amsmath,amssymb,amsfonts}
\usepackage{algorithmic}
\usepackage{graphicx}
\usepackage{textcomp}
\usepackage{xcolor}
\def\BibTeX{{\rm B\kern-.05em{\sc i\kern-.025em b}\kern-.08em
    T\kern-.1667em\lower.7ex\hbox{E}\kern-.125emX}}
\begin{document}

\title{Paper Review: Security Practices and Challenges in Game Development\\

\thanks{Identify applicable funding agency here. If none, delete this.}
}

\author{\IEEEauthorblockN{1\textsuperscript{st} Joseth S. Valdez}
\IEEEauthorblockA{\textit{Department of Computer Science} \\
\textit{The University of Texas at El Paso}\\
El Paso, USA \\
jsvaldez2@miners.utep.edu}
\and
\IEEEauthorblockN{2\textsuperscript{nd} Bryan Adame}
\IEEEauthorblockA{\textit{Department of Computer Science} \\
\textit{The University of Texas at El Paso}\\
El Paso, USA \\
badame2@miners.utep.edu}
\and
\IEEEauthorblockN{3\textsuperscript{rd} Daniel Del Villar}
\IEEEauthorblockA{\textit{Department of Computer Science} \\
\textit{The University of Texas at El Paso}\\
El Paso, USA \\
ddelvillar@miners.utep.edu}
\and
\IEEEauthorblockN{4\textsuperscript{th} Alexis Garcia Enriquez}
\IEEEauthorblockA{\textit{Department of Computer Science} \\
\textit{The University of Texas at El Paso}\\
El Paso, USA \\
agarciaenr@miners.utep.edu}
\and
\IEEEauthorblockN{5\textsuperscript{th} Henry Ng}
\IEEEauthorblockA{\textit{Department of Computer Science} \\
\textit{The University of Texas at El Paso}\\
El Paso, USA \\
hng@miners.utep.edu}
\and
\IEEEauthorblockN{6\textsuperscript{th} Paola L. Renteria}
\IEEEauthorblockA{\textit{Department of Computer Science} \\
\textit{The University of Texas at El Paso}\\
El Paso, USA \\
plrenteria2@miners.utep.edu}
\linebreakand
\IEEEauthorblockN{7\textsuperscript{th} Kevin D. Llamas}
\IEEEauthorblockA{\textit{Department of Computer Science} \\
\textit{The University of Texas at El Paso}\\
El Paso, USA \\
kdllamas@miners.utep.edu}
}

\maketitle

\section{Introduction}
In this paper, we’ll review the findings of a qualitative study on security in game development directed by a team of security researchers at the CISPA Helmholtz Center for Information Security. Although the video game industry is one of the largest software markets, it has traditionally focused mainly on the entertainment aspect of games, leaving research on their security practices scarce. The authors of the study attempted to address this gap by performing a series of semi-structured interviews with industry experts to investigate their security awareness, priorities, challenges, and practices. The results of their study showed that game industry professionals are typically aware about the risks of security flaws, but they are typically constrained by limited resources such as time and budget, so they tend to treat security as a secondary priority. In this paper, we attempt to summarize the study’s main findings based on its methodology and results, and how it could be improved in future research.

\section{Key Points}


\section{Critical Analysis}


\section{Novelty}
The paper’s novelty lies in exploring the human and organizational factors that influence cybersecurity practices in game development. Whereas previous research often examined specific technical security issues, such as cheating or network attacks, this study gathers firsthand accounts through 20 semi-structured interviews with industry stakeholders, including developers, freelancers, managers, and publishers. This approach reveals different perspectives on cybersecurity priorities, awareness, and practices within game development. It explores the real-world challenges companies face, such as time and budget constraints, gaps in security knowledge, and legal requirements. The contribution extends beyond identifying technical weaknesses and threats: it helps explain why security can receive low priority even when its importance is recognized.

\section*{Acknowledgment}


\section*{References}


\end{document}
